\section{Gestión de transacciones y concurrencia}
\label{sec:transacciones}

\subsection{Unidad de trabajo}

Un caso de uso abre una transacción y la confirma una sola vez. Ese
\codigo{COMMIT} cubre, para una transferencia, los asientos del libro mayor,
los dos saldos materializados, la fila de \codigo{transfers} y el registro de
idempotencia. La figura~\ref{fig:secuencia} muestra el flujo completo.

\begin{figure}[htbp]
\centering
\resizebox{\textwidth}{!}{% Secuencia de una transferencia, con el camino de idempotencia.
\begin{tikzpicture}[
  font=\scriptsize,
  >={Stealth[length=2mm]},
  actor/.style={draw=acento!55, fill=acento!10, rounded corners=2pt,
                minimum width=22mm, minimum height=6mm,
                font=\scriptsize\bfseries},
  vida/.style={draw=tinta!35, dashed},
  msj/.style={->, draw=tinta!75},
  ret/.style={->, draw=tinta!45, dashed},
  rotulo/.style={font=\scriptsize, align=center, fill=white, inner sep=1.5pt},
  nota/.style={draw=aviso!60, fill=yellow!10, rounded corners=2pt,
               text width=54mm, inner sep=4pt, font=\scriptsize, align=left},
]

\def\xC{0}   \def\xA{4.2}  \def\xS{8.6}  \def\xD{13.2}
\def\yBot{-12.2}

\node[actor] (C) at (\xC,0) {Cliente};
\node[actor] (A) at (\xA,0) {API};
\node[actor] (S) at (\xS,0) {Servicio};
\node[actor] (D) at (\xD,0) {PostgreSQL};

\foreach \x in {\xC,\xA,\xS,\xD} \draw[vida] (\x,-0.35) -- (\x,\yBot);

% Barra de actividad de la transacción
\fill[acento!20, draw=acento!45] (\xD-0.09,-2.5) rectangle (\xD+0.09,-10.9);

% --- Mensajes ---
\draw[msj] (\xC,-1.0) -- (\xA,-1.0);
\node[rotulo, above] at ({(\xC+\xA)/2},-1.0)
  {\codigo{POST /transfers}\\[-2pt]{\tiny con \codigo{Idempotency-Key}}};

\draw[msj] (\xA,-1.9) -- (\xS,-1.9);
\node[rotulo, above] at ({(\xA+\xS)/2},-1.9)
  {valida esquema y token;\\[-2pt]{\tiny la cuenta origen sale del \codigo{sub}}};

\draw[msj] (\xS,-2.7) -- (\xD,-2.7);
\node[rotulo, above] at ({(\xS+\xD)/2},-2.7) {\codigo{BEGIN}};

\draw[msj] (\xS,-3.5) -- (\xD,-3.5);
\node[rotulo, above] at ({(\xS+\xD)/2},-3.5)
  {\codigo{INSERT INTO idempotency\_keys}\\[-2pt]\codigo{ON CONFLICT DO NOTHING}};

\draw[ret] (\xD,-4.2) -- (\xS,-4.2);
\node[rotulo, above] at ({(\xS+\xD)/2},-4.2) {clave reclamada};

\draw[msj] (\xS,-5.1) -- (\xD,-5.1);
\node[rotulo, above] at ({(\xS+\xD)/2},-5.1)
  {\codigo{SELECT ... id=:a FOR UPDATE}\\[-2pt]
   \codigo{SELECT ... id=:b FOR UPDATE}};

\node[nota, anchor=north west] at (\xC-0.4,-5.9) {%
  \textbf{Orden ascendente de identificador.} Dos sentencias separadas y no un
  \codigo{IN (...) ORDER BY}: el orden en que el motor bloquea un conjunto
  depende del plan de ejecución. Así la espera circular es imposible.};

\draw[msj] (\xS,-6.6) -- (\xD,-6.6);
\node[rotulo, above] at ({(\xS+\xD)/2},-6.6)
  {\codigo{INSERT INTO transfers}};

\draw[msj] (\xS,-7.4) -- (\xD,-7.4);
\node[rotulo, above] at ({(\xS+\xD)/2},-7.4)
  {\codigo{INSERT} dos asientos ($-x$, $+x$)\\[-2pt]
   \codigo{UPDATE} los dos saldos};

\draw[msj] (\xS,-8.4) -- (\xD,-8.4);
\node[rotulo, above] at ({(\xS+\xD)/2},-8.4)
  {\codigo{UPDATE idempotency\_keys}\\[-2pt]{\tiny con la respuesta a reproducir}};

\draw[msj] (\xS,-9.3) -- (\xD,-9.3);
\node[rotulo, above] at ({(\xS+\xD)/2},-9.3) {\codigo{COMMIT}};

\node[nota, anchor=north west, draw=acento!60, fill=acento!8]
  at (\xC-0.4,-9.9) {%
  \textbf{Un único \codigo{COMMIT}} confirma asientos, saldos y registro de
  idempotencia. No hay ningún instante en el que el dinero se haya movido y la
  clave no esté escrita.};

\draw[ret] (\xA,-11.8) -- (\xC,-11.8);
\node[rotulo, above] at ({(\xC+\xA)/2},-11.8) {\codigo{201 Created}};

\end{tikzpicture}
}
\caption{Secuencia de una transferencia.}
\label{fig:secuencia}
\end{figure}

\subsection{Dos estrategias implementadas}

Ambas se seleccionan con la variable de entorno
\codigo{PAYCORE\_CONCURRENCY\_STRATEGY}. No se trata de tener dos caminos en
producción, sino de poder publicar una comparación con datos propios en lugar
de una preferencia argumentada.

\paragraph{\codigo{pessimistic} (por defecto).} Nivel Read Committed más
\codigo{SELECT ... FOR UPDATE} sobre las filas de las dos cuentas, adquiridas
siempre \textbf{en orden ascendente de identificador}.

\paragraph{\codigo{serializable}.} Nivel \codigo{SERIALIZABLE} sin bloqueos de
fila, con reintento ante \sqlstate{SQLSTATE 40001}, hasta tres intentos, con
retardo exponencial y desviación aleatoria. La documentación de PostgreSQL
indica expresamente que una aplicación que use este nivel debe disponer de una
forma generalizada de manejar los fallos de serialización~\cite{pg16iso}.

Que este modo \emph{no} tome bloqueos de fila es deliberado: con ellos las dos
estrategias dejarían de ser distinguibles y la comparación no mediría nada.

\subsection{Por qué dos sentencias y no un \codigo{IN (...) ORDER BY}}

La primera versión del bloqueo usaba
\codigo{WHERE id IN (:a, :b) ORDER BY id FOR UPDATE}. Es una formulación común
y es \emph{insuficiente}: el orden en que PostgreSQL adquiere los bloqueos de
un conjunto depende del plan de ejecución y no está garantizado por la
cláusula de ordenación. La implementación final emite \textbf{dos sentencias
separadas} en orden ascendente, de modo que el orden lo impone la aplicación.
Como todas las instancias ordenan igual, la espera circular es imposible.

La documentación del motor recomienda exactamente esta defensa: adquirir los
bloqueos sobre múltiples objetos en un orden consistente~\cite{pg16lock}.

\subsection{El deadlock no se reintenta}

La función que envuelve la transacción reintenta ante \sqlstate{40001} pero
\textbf{no} ante \sqlstate{40P01}, que es el código de deadlock. Es una
decisión deliberada: el orden determinista de bloqueo debe hacerlo imposible,
y reintentarlo enmascararía un fallo de diseño en lugar de exponerlo. Una
prueba lanza cincuenta transferencias cruzadas simultáneas, que es el ejemplo
de libro de la espera circular, y exige que las cincuenta se completen sin un
solo error. \estado{verificado}

\subsection{Un defecto real: la actualización perdida del mapa de identidad}

\label{subsec:lostupdate}
La prueba de concurrencia con cien peticiones aprobó \textbf{88 transferencias
de 100} cuando el máximo posible eran 3. Las 76 pruebas existentes en ese
momento pasaban todas.

El bloqueo se adquiría correctamente. El defecto estaba una capa más arriba.
El caso de uso leía la cuenta de origen \emph{sin} bloqueo antes de volver a
leerla \emph{con} bloqueo. Cuando una entidad ya está en el mapa de identidad
de la sesión de SQLAlchemy, la biblioteca devuelve ese mismo objeto sin
refrescar sus atributos: la sentencia \codigo{FOR UPDATE} se ejecutaba y el
bloqueo se tomaba, pero la comprobación de fondos evaluaba el saldo obsoleto
leído antes del bloqueo. Eso es exactamente una actualización perdida.

La corrección es una línea:

\begin{lstlisting}[language=Python,caption={Sin \codigo{populate\_existing} el
bloqueo se toma pero se razona sobre datos antiguos.},label={lst:populate}]
cuenta = session.execute(
    consulta.execution_options(populate_existing=True)
).scalar_one_or_none()
\end{lstlisting}

Tres lecciones que este defecto deja documentadas:

\begin{enumerate}[leftmargin=1.4em]
  \item Ninguna prueba de un solo hilo podía detectarlo.
  \item Tomar el bloqueo correcto en SQL no basta si la capa de mapeo
    devuelve datos cacheados.
  \item La restricción \codigo{CHECK} de no negatividad no llegó a dispararse,
    porque el saldo se calculaba a partir del valor obsoleto y la fila escrita
    nunca cruzaba el cero. Una restricción de motor acota el daño, pero no
    sustituye a la prueba.
\end{enumerate}
